\begin{figure}[t]\centering
\begin{tikzpicture}[x=1cm,y=1cm,line width=0.4pt,font=\footnotesize]
  \draw[fill=black!8,draw=black!40] (-0.55,1.62) rectangle (0.55,2.22);
  \draw[fill=black!8,draw=black!40] (1.70,-0.30) rectangle (2.55,0.30);
  \draw[fill=black!8,draw=black!40] (-2.40,-0.30) rectangle (-1.60,0.30);
  \draw[densely dotted,black!40] (0,0) circle (1.12);
  \foreach \a in {0,22.5,45,67.5,112.5,135,157.5,180,202.5,225,247.5,270,292.5,315,337.5}{
    \draw[-{Stealth[length=1.6mm]},black!40] (\a:0.64) -- (\a:1.08);}
  \draw[-{Stealth[length=2.4mm]},black,line width=1pt] (90:0.64) -- (90:1.08);
  \draw[fill=orange!30,draw=orange!80!black,line width=0.8pt]
       (-0.5,-0.28) rectangle (0.5,0.28);
  \node[font=\scriptsize,black!65] at (33.75:1.38) {$d$};
  \node[anchor=west,font=\scriptsize,align=left] at (0.75,1.30)
       {a move that lowers the criterion:\\the box is not held};
\end{tikzpicture}
\caption{The test, one box and one criterion. The box is moved a distance $d$, two per
cent of the drawing's width, in each of sixteen exact directions, and the criterion is
recomputed. If some move lowers it (the dark arrow), the box is not held; if every move
raises it, the box is \emph{held}.}
\label{fig:test}
\end{figure}
